\ifdefined\gkver\else\def\gkver{student}\fi
\documentclass[\gkver]{gaokaozhenti}
\begin{document}

\chapter{2021年广东}

\begin{enumerate}
\item
%% number: 1
%% paperName: 2021年广东高考第 1 题 4 分
%% typeId: 1
%% type: 单选题
%% chapter: 原子和原子核
%% point: 半衰期
%% method: 无
%% score: 4
%% degree: 850
%% duplicateId: 0
%% body:
科学家发现银河系中存在大量的放射性同位素铝26，铝26的半衰期为72万年，其衰变方程为 \ce{^{26}_{13}Al -> ^{26}_{12}Mg + Y}，下列说法正确的是 \xzanswer{C}

\fourchoices[answer=C]
{Y 是氦核}
{Y 是质子}
{再经过 72 万年，现有的铝 26 衰变一半}
{再经过 144 万年，现有的铝 26 全部衰变}

%% answer: C
%% memo:
\memoanswer{
	AB．根据核反应的质量数和电荷数守恒可知，该核反应是

	\ce{^{26}_{13}Al -> ^{26}_{12}Mg + ^{0}_{1}e}

	即 Y 是正电子，选项 AB 错误；

	CD．因 72 万年是一个半衰期，可知再过 72 万年，现有的铝 26 衰变一半；再过 144 万年，即两个半衰期，现有的铝 26 衰变四分之三，选项 C 正确，D 错误；

	故选 C。
}

\item
%% number: 2
%% paperName: 2021年广东高考第 2 题 4 分
%% typeId: 1
%% type: 单选题
%% chapter: 曲线运动
%% point: 人造地球卫星
%% method: 无
%% score: 4
%% degree: 650
%% duplicateId: 0
%% body:
2021年4月，我国自主研发的空间站“天和”核心舱成功发射并入轨运行，若核心舱绕地球的运行可视为匀速圆周运动，已知引力常量，由下列物理量能计算出地球质量的是 \xzanswer{D}

\fourchoices[answer=D]
{核心舱的质量和绕地半径}
{核心舱的质量和绕地周期}
{核心舱的绕地角速度和绕地周期}
{核心舱的绕地线速度和绕地半径}

%% answer: D
%% memo:
\memoanswer{
	根据核心舱做圆周运动的向心力由地球的万有引力提供，可得

	$G\frac{Mm}{r^{2}}=m\frac{v^{2}}{r}=m\omega^{2}r=m\frac{4\pi^{2}}{T^{2}}r$

	可得

	$M=\frac{v^{2}r}{G}=\frac{\omega^{2}r^{3}}{G}=\frac{4\pi^{2}r^{3}}{GT^{2}}$

	可知已知核心舱的质量和绕地半径、已知核心舱的质量和绕地周期以及已知核心舱的角速度和绕地周期，都不能求解地球的质量；若已知核心舱的绕地线速度和绕地半径可求解地球的质量。

	故选 D。
}

\item
%% number: 3
%% paperName: 2021年广东高考第 3 题 4 分
%% typeId: 1
%% type: 单选题
%% chapter: 运动和力
%% point: 牛顿第二定律
%% method: 无
%% score: 4
%% degree: 650
%% duplicateId: 0
%% body:
唐代《来耜经》记载了曲辕犁相对直辕犁的优势之一是起土省力，设牛用大小相等的拉力 $F$ 通过耕索分别拉两种犁，$F$ 与竖直方向的夹角分别为 $\alpha$ 和 $\beta$，$\alpha<\beta$，如图所示，忽略耕索质量，耕地过程中，下列说法正确的是 \xzanswer{B}

\onepicture[w=8cm]{figs/03.png}

\fourchoices[answer=B]
{耕索对曲辕犁拉力的水平分力比对直辕犁的大}
{耕索对曲辕犁拉力的竖直分力比对直辕犁的大}
{曲辕犁匀速前进时，耕索对犁的拉力小于犁对耕索的拉力}
{直辕犁加速前进时，耕索对犁的拉力大于犁对耕索的拉力}

%% answer: B
%% memo:
\memoanswer{
	A．将拉力 $F$ 正交分解如图所示

	\onepicture[w=4.5cm]{figs/03_an.png}

	则在 $x$ 方向可得出：$F_{x\text{曲}}=F\sin\alpha$，$F_{x\text{直}}=F\sin\beta$

	在 $y$ 方向可得出：$F_{y\text{曲}}=F\cos\alpha$，$F_{y\text{直}}=F\cos\beta$

	由题知 $\alpha<\beta$，则 $\sin\alpha<\sin\beta$，$\cos\alpha>\cos\beta$

	则可得到：$F_{x\text{曲}}<F_{x\text{直}}$，$F_{y\text{曲}}>F_{y\text{直}}$

	A 错误、B 正确；

	CD．耕索对犁的拉力与犁对耕索的拉力是一对相互作用力，它们大小相等，方向相反，无论是加速还是匀速，则 CD 错误。

	故选 B。
}

\item
%% number: 4
%% paperName: 2021年广东高考第 4 题 4 分
%% typeId: 1
%% type: 单选题
%% chapter: 曲线运动
%% point: 线速度角速度
%% method: 无
%% score: 4
%% degree: 550
%% duplicateId: 0
%% body:
由于高度限制，车库出入口采用图所示的曲杆道闸，道闸由转动杆 OP 与横杆 PQ 链接而成，P、Q 为横杆的两个端点。在道闸抬起过程中，杆 PQ 始终保持水平。杆 OP 绕 O 点从与水平方向成 $30\Udu$ 匀速转动到 $60\Udu$ 的过程中，下列说法正确的是 \xzanswer{A}

\onepicture[w=5cm]{figs/04.png}

\fourchoices[answer=A]
{P 点的线速度大小不变}
{P 点的加速度方向不变}
{Q 点在竖直方向做匀速运动}
{Q 点在水平方向做匀速运动}

%% answer: A
%% memo:
\memoanswer{
	A．由题知杆 OP 绕 O 点从与水平方向成 $30\Udu$ 匀速转动到 $60\Udu$，则 P 点绕 O 点做匀速圆周运动，则 P 点的线速度大小不变，A 正确；

	B．由题知杆 OP 绕 O 点从与水平方向成 $30\Udu$ 匀速转动到 $60\Udu$，则 P 点绕 O 点做匀速圆周运动，P 点的加速度方向时刻指向 O 点，B 错误；

	C．Q 点在竖直方向的运动与 P 点相同，位移 $y$ 关于时间 $t$ 的关系为：$y=l_{OP}\sin\left(\frac{\pi}{6}+\omega t\right)$

	则可看出 Q 点在竖直方向不是匀速运动，C 错误；

	D．Q 点在水平方向的位移 $x$ 关于时间 $t$ 的关系为

	$x=l_{OP}\cos\left(\frac{\pi}{6}+\omega t\right)+l_{PQ}$

	则可看出 Q 点在水平方向也不是匀速运动，D 错误。

	故选 A。
}

\item
%% number: 5
%% paperName: 2021年广东高考第 5 题 4 分
%% typeId: 1
%% type: 单选题
%% chapter: 磁场
%% point: 左手定则
%% method: 无
%% score: 4
%% degree: 550
%% duplicateId: 0
%% body:
截面为正方形的绝缘弹性长管中心有一固定长直导线，长管外表面固定着对称分布的四根平行长直导线，若中心直导线通入电流 $I_{1}$，四根平行直导线均通入电流 $I_{2}$，$I_{1}\gg I_{2}$，电流方向如图所示，下列截面图中可能正确表示通电后长管发生形变的是 \xzanswer{C}

\onepicture[w=3.5cm]{figs/05.png}

\fourchoices[answer=C, ispicture=true, h=2.6cm]
{figs/05a.png}
{figs/05b.png}
{figs/05c.png}
{figs/05d.png}

%% answer: C
%% memo:
\memoanswer{
	因 $I_{1}\gg I_{2}$，则可不考虑四个边上的直导线之间的相互作用；根据两通电直导线间的安培力作用满足“同向电流相互吸引，异向电流相互排斥”，则正方形左右两侧的直导线 $I_{2}$ 要受到 $I_{1}$ 吸引的安培力，形成凹形，正方形上下两边的直导线 $I_{2}$ 要受到 $I_{1}$ 排斥的安培力，形成凸形，故变形后的形状如图 C。

	故选 C。
}

\item
%% number: 6
%% paperName: 2021年广东高考第 6 题 4 分
%% typeId: 1
%% type: 单选题
%% chapter: 电场
%% point: 等势面
%% method: 无
%% score: 4
%% degree: 550
%% duplicateId: 0
%% body:
图是某种静电推进装置的原理图，发射极与吸极接在高压电源两端，两极间产生强电场，虚线为等势面，在强电场作用下，一带电液滴从发射极加速飞向吸极，a、b 是其路径上的两点，不计液滴重力，下列说法正确的是 \xzanswer{D}

\onepicture[w=5.5cm]{figs/06.png}

\fourchoices[answer=D]
{a 点的电势比 b 点的低}
{a 点的电场强度比 b 点的小}
{液滴在 a 点的加速度比在 b 点的小}
{液滴在 a 点的电势能比在 b 点的大}

%% answer: D
%% memo:
\memoanswer{
	A．高压电源左为正极，则所加强电场的场强向右，而沿着电场线电势逐渐降低，可知 $\varphi_{a}>\varphi_{b}$，故 A 错误；

	B．等差等势线的疏密反映场强的大小，由图可知 $a$ 处的等势线较密，则 $E_{a}>E_{b}$，故 B 错误；

	C．液滴的重力不计，根据牛顿第二定律可知，液滴的加速度为 $a=\frac{qE}{m}$；因 $E_{a}>E_{b}$，可得 $a_{a}>a_{b}$，故 C 错误；

	D．液滴在电场力作用下向右加速，则电场力做正功，动能增大，电势能减少，即 $E_{pa}>E_{pb}$，故 D 正确；

	故选 D。
}

\item
%% number: 7
%% paperName: 2021年广东高考第 7 题 4 分
%% typeId: 1
%% type: 单选题
%% chapter: 交流电
%% point: 变压器
%% method: 无
%% score: 4
%% degree: 550
%% duplicateId: 0
%% body:
某同学设计了一个充电装置，如图所示，假设永磁铁的往复运动在螺线管中产生近似正弦式交流电，周期为 $0.2\Us$，电压最大值为 $0.05\UV$，理想变压器原线圈接螺线管，副线圈接充电电路，原、副线圈匝数比为 $1:60$，下列说法正确的是 \xzanswer{B}

\onepicture[w=7cm]{figs/07.png}

\fourchoices[answer=B]
{交流电的频率为 $10\UHz$}
{副线圈两端电压最大值为 $3\UV$}
{变压器输入电压与永磁铁磁场强弱无关}
{充电电路的输入功率大于变压器的输入功率}

%% answer: B
%% memo:
\memoanswer{
	A．周期是 $T=0.2\Us$，频率是

	$f=\frac{1}{T}=5\UHz$

	故 A 错误；

	B．由理想变压器原理可知

	$\frac{U_{1}}{U_{2}}=\frac{n_{1}}{n_{2}}$

	解得，副线圈两端的最大电压为

	$U_{2}=\frac{n_{2}}{n_{1}}U_{1}=3\UV$

	故 B 正确；

	C．根据法拉第电磁感应定律可知，永磁铁磁场强，线圈中产生的感应电动势越大，变压器的输入电压会越大，故 C 错误；

	D．由理想变压器原理可知，充电电路的输入功率等于变压器的输入功率，故 D 错误。

	故选 B。
}

\item
%% number: 8
%% paperName: 2021年广东高考第 8 题 6 分
%% typeId: 2
%% type: 多选题
%% chapter: 直线运动
%% point: 运动图象
%% method: 无
%% score: 6
%% degree: 650
%% duplicateId: 0
%% body:
赛龙舟是端午节的传统活动。下列 $v-t$ 和 $s-t$ 图像描述了五条相同的龙舟从同一起点线同时出发、沿长直河道划向同一终点线的运动全过程，其中能反映龙舟甲与其它龙舟在途中出现船头并齐的有 \xzanswer{BD}

\fourchoices[answer=BD, ispicture=true, h=2.4cm]
{figs/08a.png}
{figs/08b.png}
{figs/08c.png}
{figs/08d.png}

%% answer: BD
%% memo:
\memoanswer{
	A．此图是速度图像，由图可知，甲的速度一直大于乙的速度，所以中途不可能出现甲乙船头并齐，故 A 错误；

	B．此图是速度图像，由图可知，开始丙的速度大，后来甲的速度大，速度图像中图像与横轴围成的面积表示位移，由图可以判断在中途甲、丙位移会相同，所以在中途甲丙船头会并齐，故 B 正确；

	C．此图是位移图像，由图可知，丁一直运动在甲的前面，所以中途不可能出现甲丁船头并齐，故 C 错误；

	D．此图是位移图像，交点表示相遇，所以甲戊在中途船头会齐，故 D 正确。

	故选 BD。
}

\item
%% number: 9
%% paperName: 2021年广东高考第 9 题 6 分
%% typeId: 2
%% type: 多选题
%% chapter: 曲线运动
%% point: 平抛运动
%% method: 无
%% score: 6
%% degree: 650
%% duplicateId: 0
%% body:
长征途中，为了突破敌方关隘，战士爬上陡销的山头，居高临下向敌方工事内投掷手榴弹，战士在同一位置先后投出甲、乙两颗质量均为 $m$ 的手榴弹，手榴弹从投出的位置到落地点的高度差为 $h$，在空中的运动可视为平抛运动，轨迹如图所示，重力加速度为 $g$，下列说法正确的有 \xzanswer{BC}

\onepicture[w=4.5cm]{figs/09.png}

\fourchoices[answer=BC]
{甲在空中的运动时间比乙的长}
{两手榴弹在落地前瞬间，重力的功率相等}
{从投出到落地，每颗手榴弹的重力势能减少 $mgh$}
{从投出到落地，每颗手榴弹的机械能变化量为 $mgh$}

%% answer: BC
%% memo:
\memoanswer{
	A．由平抛运动规律可知，做平抛运动的时间

	$t=\sqrt{\frac{2h}{g}}$

	因为两手榴弹运动的高度差相同，所以在空中运动时间相等，故 A 错误；

	B．做平抛运动的物体落地前瞬间重力的功率

	$P=mgv\cos\theta=mgv_{y}=mg\sqrt{2gh}$

	因为两手榴弹运动的高度差相同，质量相同，所以落地前瞬间，两手榴弹重力功率相同，故 B 正确；

	C．从投出到落地，手榴弹下降的高度为 $h$，所以手榴弹重力势能减小量

	$\Delta E_{p}=mgh$

	故 C 正确；

	D．从投出到落地，手榴弹做平抛运动，只有重力做功，机械能守恒，故 D 错误。

	故选 BC。
}

\item
%% number: 10
%% paperName: 2021年广东高考第 10 题 6 分
%% typeId: 2
%% type: 多选题
%% chapter: 电磁感应
%% point: 切割磁感线电动势
%% method: 无
%% score: 6
%% degree: 550
%% duplicateId: 0
%% body:
如图所示，水平放置足够长光滑金属导轨 $abc$ 和 $de$，$ab$ 与 $de$ 平行，$bc$ 是以 $O$ 为圆心的圆弧导轨，圆弧 $be$ 左侧和扇形 $Obc$ 内有方向如图的匀强磁场，金属杆 $OP$ 的 $O$ 端与 $e$ 点用导线相接，$P$ 端与圆弧 $bc$ 接触良好，初始时，可滑动的金属杆 $MN$ 静止在平行导轨上，若杆 $OP$ 绕 $O$ 点在匀强磁场区内从 $b$ 到 $c$ 匀速转动时，回路中始终有电流，则此过程中，下列说法正确的有 \xzanswer{AD}

\onepicture[w=6cm]{figs/10.png}

\fourchoices[answer=AD]
{杆 OP 产生的感应电动势恒定}
{杆 OP 受到的安培力不变}
{杆 MN 做匀加速直线运动}
{杆 MN 中的电流逐渐减小}

%% answer: AD
%% memo:
\memoanswer{
	A．OP 转动切割磁感线产生的感应电动势为

	$E=\frac{1}{2}Br^{2}\omega$

	因为 OP 匀速转动，所以杆 OP 产生的感应电动势恒定，故 A 正确；

	BCD．杆 OP 匀速转动产生的感应电动势产生的感应电流由 M 到 N 通过 MN 棒，由左手定则可知，MN 棒会向左运动，MN 棒运动会切割磁感线，产生电动势与原来电流方向相反，让回路电流减小，MN 棒所受合力为安培力，电流减小，安培力会减小，加速度减小，故 D 正确，BC 错误。

	故选 AD。
}

\item
%% number: 11
%% paperName: 2021年广东高考第 11 题 7 分
%% typeId: 4
%% type: 实验题
%% chapter: 力物体的平衡
%% point: 胡克定律
%% method: 无
%% score: 7
%% degree: 650
%% duplicateId: 0
%% body:
某兴趣小组测量一缓冲装置中弹簧的劲度系数，缓冲装置如图所示，固定在斜面上的透明有机玻璃管与水平面夹角为 $30\Udu$，弹簧固定在有机玻璃管底端。实验过程如下：先沿管轴线方向固定一毫米刻度尺，再将单个质量为 $200\Ug$ 的钢球（直径略小于玻璃管内径）逐个从管口滑进，每滑进一个钢球，待弹簧静止，记录管内钢球的个数 $n$ 和弹簧上端对应的刻度尺示数 $L_{0}$，数据如表所示。实验过程中弹簧始终处于弹性限度内。采用逐差法计算弹簧压缩量，进而计算其劲度系数。

\begin{table}[htbp]
\centering
\begin{tblr}{|c|c|c|c|c|c|c|}
\hline
$n$ & 1 & 2 & 3 & 4 & 5 & 6 \\
\hline
$L_{n}/\Ucm[a]$ & 8.04 & 10.03 & 12.05 & 14.07 & 16.11 & 18.09 \\
\hline
\end{tblr}
\end{table}

\begin{enumerate}
	\item
	利用 $\Delta L_{i}=L_{i+3}-L_{i}$（$i=1$，$2$，$3$）计算弹簧的压缩量：$\Delta L_{1}=6.03\Ucm$，$\Delta L_{2}=6.08\Ucm$，$\Delta L_{3}=$\tkanswer{6.04}$\Ucm$，压缩量的平均值 $\overline{\Delta L}=\frac{\Delta L_{1}+\Delta L_{2}+\Delta L_{3}}{3}=$\tkanswer{6.05}$\Ucm$；

	\item
	上述 $\overline{\Delta L}$ 是管中增加\tkanswer{3}个钢球时产生的弹簧平均压缩量；

	\item
	忽略摩擦，重力加速度 $g$ 取 $9.80\Umsq$，该弹簧的劲度系数为\tkanswer{48.6}$\UNm$。（结果保留 3 位有效数字）
\end{enumerate}

\onepicture[w=4.5cm, align=right]{figs/11.png}

%% answer:
\jdanswer{
\begin{enumerate}
	\item
	6.04，6.05

	\item
	3

	\item
	48.6
\end{enumerate}
}
%% memo:
\memoanswer{
	（1）根据压缩量的变化量为

	$\Delta L_{3}=L_{6}-L_{3}=(18.09-12.05)\Ucm=6.04\Ucm$

	压缩量的平均值为

	$\overline{\Delta L}=\frac{\Delta L_{1}+\Delta L_{2}+\Delta L_{3}}{3}=\frac{6.03+6.08+6.04}{3}\Ucm\approx6.05\Ucm$

	（2）因三个 $\Delta L$ 是相差 3 个钢球的压缩量之差，则所求平均值为管中增加 3 个钢球时产生的弹簧平均压缩量；

	（3）根据钢球的平衡条件有

	$3mg\sin\theta=k\cdot\overline{\Delta L}$

	解得

	$k=\frac{3mg\sin\theta}{\overline{\Delta L}}=\frac{3\times0.2\times9.8\times\sin30\Udu}{6.05\times10^{-2}}\UNm\approx48.6\UNm$
}

\item
%% number: 12
%% paperName: 2021年广东高考第 12 题 9 分
%% typeId: 4
%% type: 实验题
%% chapter: 电路
%% point: 电路实验
%% method: 无
%% score: 9
%% degree: 550
%% duplicateId: 0
%% body:
某小组研究热敏电阻阻值随温度的变化规律。根据实验需要已选用了规格和量程合适的器材。

\begin{enumerate}
	\item
	先用多用电表预判热敏电阻阻值随温度的变化趋势。选择适当倍率的欧姆挡，将两表笔\tkanswer{短接}，调节欧姆调零旋钮，使指针指向右边“$0\UO$”处。测量时观察到热敏电阻温度越高，相同倍率下多用电表指针向右偏转角度越大，由此可判断热敏电阻阻值随温度的升高而\tkanswer{减小}。

	\item
	再按图连接好电路进行测量。

	\begin{steps}
		\item
		闭合开关 S 前，将滑动变阻器 $R_{1}$ 的滑片滑到\tkanswer{b}端（选填“a”或“b”）。

		将温控室的温度设置为 $T$，电阻箱 $R_{0}$ 调为某一阻值 $R_{01}$。闭合开关 S，调节滑动变阻器 $R_{1}$，使电压表和电流表的指针偏转到某一位置。记录此时电压表和电流表的示数、$T$ 和 $R_{01}$。断开开关 S。

		再将电压表与热敏电阻 C 端间的导线改接到 D 端，闭合开关 S。反复调节 $R_{0}$ 和 $R_{1}$，使电压表和电流表的示数与上述记录的示数相同。记录此时电阻箱的阻值 $R_{02}$。断开开关 S。

		\item
		实验中记录的阻值 $R_{01}$\tkanswer{大于}$R_{02}$（选填“大于”、“小于”或“等于”）。此时热敏电阻阻值 $R_{T}=$\tkanswer{$R_{01}-R_{02}$}。
	\end{steps}
\end{enumerate}

\onepicture[w=5cm, align=right]{figs/12.png}

%% answer:
\jdanswer{
\begin{enumerate}
	\item
	短接，减小

	\item
	b，大于，$R_{01}-R_{02}$
\end{enumerate}
}
%% memo:
\memoanswer{
	（1）选择倍率适当的欧姆档，将两表笔短接；欧姆表指针向右偏转角度越大，则阻值越小，可判断热敏电阻的阻值随温度升高而减小。

	（2）①闭合开关 S 前，应将滑动变阻器 $R_{1}$ 的阻值调到最大，即将滑片滑到 b 端；

	②因两次电压表和电流表的示数相同，因为

	$R_{01}=R_{02}+R_{T}$

	即

	$R_{T}=R_{01}-R_{02}$

	可知 $R_{01}$ 大于 $R_{02}$。
}

\item
%% number: 13
%% paperName: 2021年广东高考第 13 题 11 分
%% typeId: 6
%% type: 计算题
%% chapter: 运动和力
%% point: 动量守恒定律
%% method: 无
%% score: 11
%% degree: 550
%% duplicateId: 0
%% body:
算盘是我国古老的计算工具，中心带孔的相同算珠可在算盘的固定导杆上滑动，使用前算珠需要归零，如图所示，水平放置的算盘中有甲、乙两颗算珠未在归零位置，甲靠边框 b，甲、乙相隔 $s_{1}=3.5\times10^{-2}\Um$，乙与边框 a 相隔 $s_{2}=2.0\times10^{-2}\Um$，算珠与导杆间的动摩擦因数 $\mu=0.1$。现用手指将甲以 $0.4\Ums$ 的初速度拨出，甲、乙碰撞后甲的速度大小为 $0.1\Ums$，方向不变，碰撞时间极短且不计，重力加速度 $g$ 取 $10\Umsq$。

\begin{enumerate}
	\item
	通过计算，判断乙算珠能否滑动到边框 a；

	\item
	求甲算珠从拨出到停下所需的时间。
\end{enumerate}

\onepicture[w=6cm, align=right]{figs/13.png}

%% answer:
\jdanswer{
\begin{enumerate}
	\item
	能

	\item
	$0.2\Us$
\end{enumerate}
}
%% memo:
\memoanswer{
	（1）甲乙滑动时的加速度大小均为

	$a=\mu g=1\Umsq$

	甲与乙碰前的速度 $v_{1}$，则

	$v_{1}^{2}=v_{0}^{2}-2as_{1}$

	解得 $v_{1}=0.3\Ums$

	甲乙碰撞时由动量守恒定律

	$mv_{1}=mv_{2}+mv_{3}$

	解得碰后乙的速度

	$v_{3}=0.2\Ums$

	然后乙做减速运动，当速度减为零时则

	$x=\frac{v_{3}^{2}}{2a}=\frac{0.2^{2}}{2\times1}\Um=0.02\Um=s_{2}$

	可知乙恰好能滑到边框 a；

	（2）甲与乙碰前运动的时间

	$t_{1}=\frac{v_{0}-v_{1}}{a}=\frac{0.4-0.3}{1}\Us=0.1\Us$

	碰后甲运动的时间

	$t_{2}=\frac{v_{2}}{a}=\frac{0.1}{1}\Us=0.1\Us$

	则甲运动的总时间为

	$t=t_{1}+t_{2}=0.2\Us$
}

\item
%% number: 14
%% paperName: 2021年广东高考第 14 题 15 分
%% typeId: 6
%% type: 计算题
%% chapter: 磁场
%% point: 带电粒子在磁场中的运动
%% method: 无
%% score: 15
%% degree: 350
%% duplicateId: 0
%% body:
如图是一种花瓣形电子加速器简化示意图，空间有三个同心圆 a、b、c 围成的区域，圆 a 内为无场区，圆 a 与圆 b 之间存在辐射状电场，圆 b 与圆 c 之间有三个圆心角均略小于 $90\Udu$ 的扇环形匀强磁场区Ⅰ、Ⅱ和Ⅲ。各区感应强度恒定，大小不同，方向均垂直纸面向外。电子以初动能 $E_{k0}$ 从圆 b 上 P 点沿径向进入电场，电场可以反向，保证电子每次进入电场即被全程加速，已知圆 a 与圆 b 之间电势差为 $U$，圆 b 半径为 $R$，圆 c 半径为 $\sqrt{3}R$，电子质量为 $m$，电荷量为 $e$，忽略相对论效应，取 $\tan22.5\Udu=0.4$。

\begin{enumerate}
	\item
	当 $E_{k0}=0$ 时，电子加速后均沿各磁场区边缘进入磁场，且在电场内相邻运动轨迹的夹角 $\theta$ 均为 $45\Udu$，最终从 Q 点出射，运动轨迹如图中带箭头实线所示，求Ⅰ区的磁感应强度大小、电子在Ⅰ区磁场中的运动时间及在 Q 点出射时的动能；

	\item
	已知电子只要不与Ⅰ区磁场外边界相碰，就能从出射区域出射。当 $E_{k0}=keU$ 时，要保证电子从出射区域出射，求 $k$ 的最大值。
\end{enumerate}

\onepicture[w=7cm, align=right]{figs/14.png}

%% answer:
\jdanswer{
\begin{enumerate}
	\item
	$\frac{5\sqrt{eUm}}{eR}$，$\frac{\pi R\sqrt{eUm}}{4eU}$，$8eU$

	\item
	$\frac{13}{6}$
\end{enumerate}
}
%% memo:
\memoanswer{
	（1）电子在电场中加速有

	$2eU=\frac{1}{2}mv^{2}$

	在磁场Ⅰ中，由几何关系可得

	$r=R\tan22.5\Udu=0.4R$

	$B_{1}ev=m\frac{v^{2}}{r}$

	联立解得

	$B_{1}=\frac{5\sqrt{eUm}}{eR}$

	在磁场Ⅰ中的运动周期为

	$T=\frac{2\pi r}{v}$

	由几何关系可得，电子在磁场Ⅰ中运动的圆心角为

	$\varphi=\frac{5}{4}\pi$

	在磁场Ⅰ中的运动时间为

	$t=\frac{\varphi}{2\pi}T$

	联立解得

	$t=\frac{\pi R\sqrt{meU}}{4eU}$

	从 Q 点出来的动能为

	$E_{k}=8eU$

	（2）在磁场Ⅰ中的做匀速圆周运动的最大半径为 $r_{m}$，此时圆周的轨迹与Ⅰ边界相切，由几何关系可得

	$\left(\sqrt{3}R-r_{m}\right)^{2}=R^{2}+r_{m}^{2}$

	解得

	$r_{m}=\frac{\sqrt{3}}{3}R$

	由于

	$B_{1}ev_{m}=m\frac{v_{m}^{2}}{r_{m}}$

	$2eU=\frac{1}{2}mv_{m}^{2}-keU$

	联立解得

	$k=\frac{13}{6}$
}

\item
%% number: 15
%% paperName: 2021年广东高考第 15 题 6 分
%% typeId: 3
%% type: 填空题
%% chapter: 气体性质
%% point: 玻意耳定律
%% method: 无
%% score: 6
%% degree: 750
%% duplicateId: 0
%% body:
在高空飞行的客机上某乘客喝完一瓶矿泉水后，把瓶盖拧紧。下飞机后发现矿泉水瓶变瘪了，机场地面温度与高空客舱内温度相同。由此可判断，高空客舱内的气体压强 \tkanswer{小于}（选填“大于”、“小于”或“等于”）机场地面大气压强：从高空客舱到机场地面，矿泉水瓶内气体的分子平均动能 \tkanswer{不变}（选填“变大”、“变小”或“不变”）。

%% answer:
%% memo:
\memoanswer{
	（1）机场地面温度与高空客舱温度相同，由题意知瓶内气体体积变小，以瓶内气体为研究对象，根据波意耳定律 $pV=C$，故可知高空客舱内的气体压强小于机场地面大气压强；

	（2）由于温度是平均动能的标志，气体的平均动能只与温度有关，机场地面温度与高空客舱温度相同，故从高空客舱到机场地面，瓶内气体的分子平均动能不变。
}

\item
%% number: 16
%% paperName: 2021年广东高考第 16 题 6 分
%% typeId: 6
%% type: 计算题
%% chapter: 气体性质
%% point: 玻意耳定律
%% method: 无
%% score: 6
%% degree: 650
%% duplicateId: 0
%% body:
为方便抽取密封药瓶里的药液，护士一般先用注射器注入少量气体到药瓶里后再抽取药液，如图所示，某种药瓶的容积为 $0.9\UmL$，内装有 $0.5\UmL$ 的药液，瓶内气体压强为 $1.0\times10^{5}\UPa$，护士把注射器内横截面积为 $0.3\Ucmq$、长度为 $0.4\Ucm$、压强为 $1.0\times10^{5}\UPa$ 的气体注入药瓶，若瓶内外温度相同且保持不变，气体视为理想气体，求此时药瓶内气体的压强。

\onepicture[h=4cm, align=right]{figs/16.png}

%% answer:
\jdanswer{
$1.3\times10^{5}\UPa$
}
%% memo:
\memoanswer{
	以注入后的所有气体为研究对象，由题意可知瓶内气体发生等温变化，设瓶内气体体积为 $V_{1}$，有

	$V_{1}=0.9\UmL-0.5\UmL=0.4\UmL=0.4\Ucmc$

	注射器内气体体积为 $V_{2}$，有

	$V_{2}=0.3\times0.4\Ucmc=0.12\Ucmc$

	根据理想气体状态方程有

	$p_{0}\left(V_{1}+V_{2}\right)=p_{1}V_{1}$

	代入数据解得

	$p_{1}=1.3\times10^{5}\UPa$
}

\item
%% number: 17
%% paperName: 2021年广东高考第 17 题 6 分
%% typeId: 1
%% type: 单选题
%% chapter: 机械振动
%% point: 振动图象
%% method: 无
%% score: 6
%% degree: 650
%% duplicateId: 0
%% body:
如图所示，一个轻质弹簧下端挂一小球，小球静止。现将小球向下拉动距离 $A$ 后由静止释放，并开始计时，小球在竖直方向做简谐运动，周期为 $T$。经 $\frac{T}{8}$ 时间，小球从最低点向上运动的距离 \tkanswer{小于} $\frac{A}{2}$（选填“大于”、“小于”或“等于”）；在 $\frac{T}{4}$ 时刻，小球的动能 \tkanswer{最大}（选填“最大”或“最小”）。

\onepicture[w=5.5cm]{figs/17.png}

%% answer:
%% memo:
\memoanswer{
	（1）根据简谐振动的位移公式

	$y=-A\cos\left(\frac{2\pi}{T}t\right)$

	则 $t=\frac{T}{8}$ 时有

	$y=-A\sin\left(\frac{2\pi}{T}\times\frac{T}{8}\right)=-\frac{\sqrt{2}}{2}A$

	所以小球从最低点向上运动的距离为

	$\Delta y=A-\frac{\sqrt{2}}{2}A=\frac{2-\sqrt{2}}{2}A<\frac{1}{2}A$

	则小球从最低点向上运动的距离小于 $\frac{A}{2}$。

	（2）在 $t=\frac{T}{4}$ 时，小球回到平衡位置，具有最大的振动速度，所以小球的动能最大。
}

\item
%% number: 18
%% paperName: 2021年广东高考第 18 题 6 分
%% typeId: 6
%% type: 计算题
%% chapter: 光学
%% point: 全反射
%% method: 无
%% score: 6
%% degree: 850
%% duplicateId: 0
%% body:
如图所示，一种光学传感器是通过接收器 Q 接收到光的强度变化而触发工作的。光从挡风玻璃内侧 P 点射向外侧 M 点再折射到空气中，测得入射角为 $\alpha$，折射角为 $\beta$；光从 P 点射向外侧 N 点，刚好发生全反射并被 Q 接收，求光从玻璃射向空气时临界角 $\theta$ 的正弦值表达式。

\onepicture[w=5.5cm, align=right]{figs/18.png}

%% answer:
\jdanswer{
$\frac{\sin\alpha}{\sin\beta}$
}
%% memo:
\memoanswer{
	根据光的折射定律有：$n=\frac{\sin\beta}{\sin\alpha}$

	根据光的全反射规律有：$\sin\theta=\frac{1}{n}$

	联立解得：$\sin\theta=\frac{\sin\alpha}{\sin\beta}$
}

\end{enumerate}

\end{document}
